\documentclass[10pt,a4paper]{amsart}
\usepackage[margin=19mm]{geometry}
\usepackage{amsmath,amssymb,mathtools}
\usepackage[scr=rsfs,cal=euler]{mathalfa}
\usepackage{xfrac}
\usepackage[unicode,hidelinks,bookmarks=false]{hyperref}
\newcommand{\ie}{i.e.\ }
\newcommand{\cf}{cf.\ }
\newcommand{\resp}{respectively\ }
\newcommand{\Subsection}{\subsection}
\providecommand{\nobreakdash}{\nobreak\hbox{-}}
\setlength{\emergencystretch}{2em}
\multlinegap0pt
\usepackage{bm}
\usepackage{bbm}
\usepackage{dsfont}
\usepackage{xspace}
\usepackage{cancel}
\usepackage{mathtools}
\usepackage{amsthm}
\makeatletter
\@ifundefined{theorem}{\newtheorem{theorem}{Theorem}}{}
\makeatother
\newcommand{\CalF}{\mathcal{F}}
\newcommand{\CalX}{\mathcal X}
\newcommand{\Real}{\mathbb R}
\title[The BdryMat\'ern GP: Reliable incorporation of boundary info]{The BdryMat\'ern GP: Reliable incorporation of boundary information on irregular domains for Gaussian process modeling}
\author{Liang Ding, Simon Mak, C. F. Jeff Wu}
\date{Source: arXiv:2507.09178v1 (CC BY 4.0)}
\begin{document}
\maketitle

\noindent\emph{Source and licence.} The BdryMat\'ern GP: Reliable incorporation of boundary information on irregular domains for Gaussian process modeling. Version arXiv:2507.09178v1 (CC BY 4.0), distributed under the Creative Commons Attribution 4.0 International licence, as recorded in the arXiv licence metadata for this exact version and shown on the abstract page https://arxiv.org/abs/2507.09178v1. Full rights notice: licenses/arxiv-2507.09178-CC-BY-4.0.txt.\par\smallskip

\noindent\textit{Editorial scope.} Counted unit: the Theorem whose source label is \texttt{thm:tensor} in the article (source0.tex lines 1181--1203, source label \texttt{thm:tensor}), delivered together with the article's own complete original proof of it (source0.tex lines 1205--1250, one \texttt{proof} environment of 46 lines). No mathematics is written, repaired or re-derived: statement and proof are the article's own text line for line. Required context retained uncounted: eq:boundar\_mixed\_1D (lines 1168--1171). Every definition, display and label of those lines is retained unchanged. Omitted from this package and not redistributed: 1--1167, 1172--1180, 1204--1204, 1251--2277. Nothing inside a retained excerpt is omitted or altered: every retained body is delivered exactly as the article's lines are written, so no omission receipt of any kind accompanies this package. Citations: the retained excerpts carry no citation command at all, so no bibliography entry of the article is transcribed and no reference is redistributed. Reference adaptation: none was needed and none was made; every reference inside the delivered unit resolves inside it, so no unresolved reference or citation is produced. Printed slips kept literal: no printed word, symbol, hypothesis or number of the counted statement or of its proof was altered. Layout and class: the article's own class is \texttt{imsart} with packages including amsthm, amsmath, amsfonts, amssymb, natbib, hyperref, graphicx, lipsum, epstopdf, ragged2e, mathtools, bm, euscript, fontenc; the wrapper is the lane's pinned \texttt{amsart} editorial wrapper at 10pt on A4, which restates the theorem-like environments the retained text needs and adds the article's own macro definitions that the retained text needs (\texttt{\textbackslash CalF}, \texttt{\textbackslash CalX}, \texttt{\textbackslash Real}), with extra wrapper packages (bm, bbm, dsfont, xspace, cancel, mathtools). The delivered numbering is the unit's own. Assets: no retained span contains a figure environment, a graphics-inclusion command, a PDF-page inclusion, a table or a TikZ picture; no image file is copied into this package, the package declares no asset dependency and no image file is redistributed with it. Licence: version arXiv:2507.09178v1 of this article is distributed under the Creative Commons Attribution 4.0 International licence, as recorded in the arXiv licence metadata for this exact version and shown on the abstract page https://arxiv.org/abs/2507.09178v1. A full rights notice is delivered at licenses/arxiv-2507.09178-CC-BY-4.0.txt. Characters: every retained excerpt is pure ASCII, so the delivered engine, XeLaTeX with its default Latin Modern text font, reports no missing character.
\begin{align}
\label{eq:boundar_mixed_1D}
   f(x)+cf'(x)=0, \quad c\geq 0, \quad x=0,1,
\end{align}

\begin{theorem}
\label{thm:tensor}
Consider the one-dimensional BdryMat\'ern kernel $k_{\nu,\mathcal{B},c}$ with smoothness $\nu$ for the homogeneous Robin boundaries \eqref{eq:boundar_mixed_1D} on domain $[0,1]$. This kernel has the closed-form representation:
\begin{align}
\begin{split}
k_{\nu,\mathcal{B},c}(x,x') = k_{\nu}(x,x') - k_{hp}(x,x') - k_{hp}(x',x) + k_{hh}(x,x'),
\end{split}
\label{eq:closedform}
\end{align}
\label{thm:closedform}
where $k_{\nu}(\cdot,\cdot)$ is the base Mat\'ern kernel with smoothness $\nu > 2$ and inverse length-scale $\kappa > 0$, and:
\small
\begin{align}
\begin{split}
k_{hp}(x,x') &= \frac{1/2}{\sinh(\kappa)}\left[K_1(x-1)\left(\frac{e^{\kappa x'}}{1+c\kappa}-\frac{e^{-\kappa x'}}{1-c\kappa}\right)+K_1(x)\left(\frac{e^{\kappa(1- x')}}{1-c\kappa}-\frac{e^{-\kappa(1- x')}}{1+c\kappa}\right)\right],\\
k_{hh}(x,x')&=\frac{1}{\sinh({\kappa})^{2}} \times\\
& \quad \left(\left[\cosh(\kappa)K_2(0)-K_2(1)\right]\left(\frac{e^{\kappa(1-x-x')}}{2(1-c\kappa)^2}+\frac{e^{-\kappa(1-x-x')}}{2(1+c\kappa)^2}\right) + \left[\cosh( {\kappa})K_2(1)-K_2(0)\right]\frac{\cosh( {\kappa}(x-x'))}{1-c^2\kappa^2}\right).
\end{split}
\label{eq:closedform2}
\end{align}
\normalsize
Here, $K_1(x)= k_{\nu}(x,x)-c\partial_x k_{\nu}(x,x)$ and $K_2(x)=k_{\nu}(x,x)-c^2\partial_{x x}k_{\nu}(x,x)$.
\end{theorem}

\begin{proof}
% In \eqref{eq:boundar_mixed_1D}, Dirichlet boundary condition corresponds to $c=0$ and general Robin boundary correspons to $c>0$.
% Given the one-dimensional domain $[0,1]$, the considered one-dimensional BdryMat\'ern kernel $k_{\nu,\mathcal{B},c}$ incorporates the Robin boundary condition of the form:
% % we impose the following boundary condition that combine both the Dirichlet and Robin boundary conditions:
% \begin{align}
% \label{eq:boundar_mixed_1D}
%    f(x)+cf'(x)=0,\quad x=0,1,\quad c\geq 0.
% \end{align}
Since the BdryMat\'ern kernel (with $\nu > 2$) can be represented as the convolution between Green's functions and the Mat\'ern kernel, we have:
\begin{align}
    k_{\nu,\mathcal{B}}(x,x')=&\int_{\CalX\times\CalX}G(x,s)k_{\nu-2}(s,u)G(u,x')dsdu
    =C_{\CalF^{-1}}\int_\Real\frac{f(x,\omega)\overline{f(x',\omega)}}{(\kappa^2+\omega^2)^{\nu-3/2}}d\omega.
     \label{eq:kernel_green_1d}
\end{align}
Here, the second equality follows from the spectral density representation of $k_\nu$ and switch of integral, and $f(x,\omega)=\int_{\CalX}G(x,s)e^{i\omega s}ds$ is the solution to the one-dimensional PDE:
\begin{align}
\begin{split}
   & \left(\kappa^2-\frac{\partial^2}{\partial x^2}\right)f(x)=e^{i\omega x},\\
   &f(x)+cf'(x)=0,\quad x=0,1,\quad c\geq 0.
   \end{split}
\label{eq:green_1_character}
\end{align}
The solution to \eqref{eq:green_1_character} can be represented as a particular solution $f_p$ minus a homogeneous $f_h$, i.e., $f(x,\omega)=f_p(x,\omega)-f_h(x,\omega)$, with each part satisfying:
\begin{align*}
    &\left(\kappa^2-\frac{\partial^2}{\partial x^2}
    \right)u_p(x)=e^{i\omega x},\quad \left(\kappa^2-\frac{\partial^2}{\partial x^2}\right)u_h(x)=0,\\
   & u_h(x)+cu_h'(x)=u_p(x),\quad x=0,1.
\end{align*}
Here, both $f_p$ and $f_h$ can be solved analytically as:
\begin{align}
\begin{split}
    & f_p(x)=\frac{e^{i\omega x}}{\kappa^2+\omega^2},\\
    & f_h(x)=\frac{(e^\kappa-e^{i\omega})(1+ic\omega)e^{-\kappa x}}{2(1-c\kappa)\sinh(\kappa)(\kappa^2+\omega^2)}+\frac{(e^{i\omega}-e^{-\kappa})(1+ic\omega)e^{\kappa x}}{2(1+c\kappa)\sinh(\kappa)(\kappa^2+\omega^2)}.
\end{split}
\label{eq:u_p_u_h_solution}
\end{align}
Substitute \eqref{eq:u_p_u_h_solution} into \eqref{eq:kernel_green_1d}, we have the closed form of $k_{\nu,\mathcal{B},c}$ in Equations \eqref{eq:closedform} and \eqref{eq:closedform2}.
% \begin{align}
%     & k^{\rm BM}_{\nu+2}(x,x')=k_{\nu+2}(x,x')-k_{hp}(x,x')-k_{hp}(x',x)+k_{hh}(x,x')  \label{eq:kernel_green_1d_closed} \\
%       & k_{hp}(x,x')=\frac{1/2}{\sinh(\kappa)}\left[K_1(x-1)(\frac{e^{\kappa x'}}{1+c\kappa}-\frac{e^{-\kappa x'}}{1-c\kappa})+K_1(x)(\frac{e^{\kappa(1- x')}}{1-c\kappa}-\frac{e^{-\kappa(1- x')}}{1+c\kappa})\right],\nonumber\\
%     & k_{hh}(x,x')=\bigg(\left[\cosh(\kappa)K_2(0)-K_2(1)\right]\left(\frac{e^{\kappa(1-x-x')}}{2(1-c\kappa)^2}+\frac{e^{-\kappa(1-x-x')}}{2(1+c\kappa)^2}\right)\nonumber\\
%     &\quad+ \left[\cosh( {\kappa})K_2(1)-K_2(0)\right]\frac{\cosh( {\kappa}(x-x'))}{1-c^2\kappa^2}\bigg)\frac{1}{(\sinh({\kappa}))^2},\nonumber\\
%     &K_1(x)=\int_\Real\frac{e^{i\omega x}(1-ic\omega)}{(\kappa^2+\omega^2)^{\nu+5/2}}d\omega= k_{\nu+2}(x)-c\partial_x k_{\nu+2}(x),\nonumber\\
%     &K_2(x)=\int_\Real\frac{e^{i\omega x}(1+c^2\omega^2)}{(\kappa^2+\omega^2)^{\nu+5/2}}d\omega=k_{\nu+2}(x)-c^2\partial_{x x}k_{\nu+2}(x).\nonumber
% \end{align}
\end{proof}

\end{document}
